%% =====================================================================================
%% 3 Game, benchmarks, designs.
%% Proposition 1 is checked by exhaustive search in crsd/tests/test_paper_equilibrium.py
%% (all symmetric profiles plus 6,000 asymmetric ones at 11 risk levels). If you change
%% the statement, change the test. p = 1 is excluded on purpose: at p = 1 every profile
%% that misses the target is a weak equilibrium.
%% The verbatim prompts (base, no hint, neutral wording, goal stated, two rewordings) live
%% in supplement/supplement.tex, rendered from kaggle/benchmarks/crg_task_server.py. The
%% Prompts paragraph below must quote exactly the spans each condition changes.
%% =====================================================================================
\section{The game and what rational play looks like}
\label{sec:game}

\subsection{Rules}
Six players each start with $e=40$ units. A \emph{seat} is one player position and a
\emph{table} is the six seats of one game. In each of $R=10$ rounds, every player privately
chooses a contribution $a_i^t\in\{0,2,4\}$ to a shared \emph{pool}; after each round all
players see every seat's past contributions. Let $x_i=\sum_t a_i^t$ be a player's total and
$X=\sum_i x_i$ the pool. The target is $T=120$. A player keeps what it did not pay, but if
the pool misses the target a catastrophe destroys all savings with probability $p$:
\begin{equation}
u_i \;=\;
\begin{cases}
e-x_i & \text{if } X\ge T,\\
0 \text{ with probability } p, \text{ else } e-x_i & \text{if } X<T.
\end{cases}
\label{eq:payoff}
\end{equation}
The fair share is $T/6=20$ per player, or 2 per round.

\begin{proposition}
\label{prop:eq}
Consider the game in which each player picks a total $x_i\in\{0,2,\dots,40\}$, and let
$p<1$. A profile is a Nash equilibrium if and only if $x=0$, or $X=T$ and $x_i\le pe$ for
every $i$. A target-reaching equilibrium exists if and only if $p\ge\popt=T/(6e)=\tfrac12$.
The welfare-optimal outcome is $x=0$ for $p<\popt$ and $X=T$ for $p>\popt$, so the best
attainable payoff per player is $\mathrm{opt}(p)=\max\{(1-p)e,\;e-T/6\}$.
\end{proposition}

\noindent\emph{Proof sketch.} If $X<T$, any player with $x_i>0$ gains by paying nothing, so
the only equilibrium that misses the target is $x=0$. If $X>T$, a player can pay 2 less and
still reach the target. If $X=T$, the best deviation is to pay nothing, worth $(1-p)e$
against $e-x_i$, which gives $x_i\le pe$; summing over players gives $p\ge T/(6e)$. Welfare
is $6e-X$ when the target is met and $(1-p)(6e-X)$ otherwise, so it peaks at $X=T$ or
$X=0$.\hfill$\square$

The ten-round game, whose strategies can react to the history, has the same totals in
pure-strategy Nash equilibrium; we claim nothing about subgame perfection. Paying nothing
in every round secures $(1-p)e$ whatever the others do, so no equilibrium misses the target
with $x\neq0$ or has $x_i>pe$. A payment
beyond the target can be dropped in the round the target is reached, or later, and nothing
played afterwards can lower that player's payoff. Each profile of
Proposition~\ref{prop:eq} is played by fixed schedules that ignore the history.

Three benchmarks follow and carry the rest of the paper. First, the \emph{optimum}, the
welfare optimum for risk-neutral players, switches at $\popt$, as does the best response to
five fair-share partners; it is a benchmark, not a prediction, since $x=0$ is an equilibrium
at every $p$. Second, paying is
\emph{dominated}, meaning that every unit paid lowers the seat's own payoff whatever else
happens, in two situations: at $p=0$, where both branches of \eqref{eq:payoff} are equal,
and once the outcome is \emph{settled}, meaning that the visible history already proves that
more payment cannot change it (the pool has reached the target, or cannot reach it even if
every seat pays 4 in every remaining round). Neither test depends on the agent's attitude to
risk or on its beliefs. Third, these two tests are the sharpest evidence of a default:
payment that continues where paying is dominated cannot follow from the stakes.

\begin{proposition}
\label{prop:table}
For two seats $i$ and $j$ in the same game, the difference in their expected payoffs over
the lottery is $c\,(x_j-x_i)$, with $c=1$ if the target was reached and $c=1-p$ otherwise.
\end{proposition}

\noindent Both seats face the same pool and the same lottery, so only their own payments
differ. Hence, for $p<1$, whoever pays less at a table earns more, whatever the group
achieves.

\subsection{Models, designs and measures}

We test five production models, one per provider, all at the low-cost tier of their family:
Claude Haiku 4.5 (Anthropic), Gemini 3.5 Flash-Lite (Google), GPT-5.6 Luna (OpenAI),
Qwen3-235B-A22B-Instruct (Alibaba) and Grok 4.20 in its non-reasoning mode (xAI). We call
them Claude Haiku, Gemini Flash-Lite, GPT Luna, Qwen and Grok; the supplementary material
gives the exact model identifiers. We chose this tier because agents that act at scale are
likely to run on it. Table~\ref{tab:designs} lists the designs. Every model plays the same
number of games in every design, and every cell holds \CntReps{} games.

\begin{table}[t]
\caption{Designs. A cell is one condition at one risk level: one model in self-play, one
model with one kind of script, or one pair of models with one split of seats. Each cell holds
ten games. The in-game question condition uses the unchanged prompt, so it also repeats the
self-play grid at three risk levels.}
\label{tab:designs}
\small
\setlength{\tabcolsep}{3pt}
\begin{tabular}{llr}
\toprule
Design & Risk levels $p$ & Games \\
\midrule
Self-play grid & $0, 0.1, \dots, 1$ & \CntGamesBaseline \\
In-game questions & $0.1, 0.5, 0.9$ & \CntGamesProbe \\
No hint & $0.1, 0.5, 0.9$ & \CntGamesNoCue \\
Neutral wording & $0, 0.1, 0.9$ & \CntGamesWording \\
Neutral wording, goal stated & $0, 0.1, 0.9$ & \CntGamesNeutral \\
Two rewordings & $0.1, 0.9$ & \CntGamesPara \\
Temperature zero & $0.1, 0.9$ & \CntGamesTempZero \\
One model, five scripts (4 kinds) & $0.1, 0.3, \dots, 0.9$ & \CntGamesScripted \\
Two models (10 pairs, 1 to 5 seats) & $0.1, 0.5, 0.9$ & \CntGamesMixed \\
\midrule
All & & \CntGames \\
\bottomrule
\end{tabular}
\end{table}

\paragraph{Prompts.}
Each round a seat receives one prompt: the rules, the full history of every seat's
contributions, and a request for one contribution. The supplementary material gives every
prompt verbatim. The base prompt calls the pool a ``climate account'', says that the group
``must'' reach the target, calls the loss a ``disaster'', names the game as a collective-risk
social dilemma, and never states the player's goal. It also states the equal split twice, as
an average of 2 per round and as an example in which paying 2 every round leaves 20. The
prompt controls change exactly these parts. \emph{No hint} deletes the two equal-split
phrases. \emph{Neutral wording} replaces the four normative or naming phrases with plain
ones, such as ``a shared account'' and ``the outcome depends on whether'' the pool reaches
120. \emph{Goal stated} keeps the neutral wording and adds an opening that says the player
is in a paid experiment where ``the only thing that matters to you is your own final cash
payoff''. Two \emph{rewordings} change the phrasing but keep every rule and the hint.

\paragraph{Protocol.}
The running pool is not printed, so a seat that wants it must add up the history. The
prompt is in English, no seat
is given a persona, and the temperature is 0.7 unless stated. A reply is read from its final
\texttt{CONTRIBUTION:} line; replies cut off at the output limit are retried with a larger
limit rather than guessed. The data contain \CntDecisions{} model decisions. Only
\CntParseFail{} replies could not be read, and both come from one reworded prompt.

\paragraph{Measures.}
We report the \emph{expected} payoff, the average a seat would earn if the catastrophe
lottery were run many times on the same contributions, and the share of $\mathrm{opt}(p)$ it
loses. We do not use realised payoffs. A catastrophe strikes when a uniform random draw falls
below $p$, and the draw depends only on the repetition number, so our games reuse ten draws in
every cell, all below \LotMaxDraw{}. Every missed game at $p\ge\popt$ therefore ended in
catastrophe, and realised payoffs would punish models that miss more than chance would.
Agents never see the draw, so their play is not affected. A seat plays the \emph{exact split}
if it pays 2 in all ten rounds. The unit of analysis is the game, since seats in one game
share a history and a lottery. Intervals are 95\% percentile bootstraps over games and tests
are permutation tests over games.
